\chapter{Diseño Experimental}

El presente estudio utiliza un diseño experimental factorial de efectos fijos con réplicas organizadas en bloques aleatorizados. Los factores analizados (controles de seguridad, variantes tecnológicas y niveles de carga) son seleccionados intencionalmente. La estructura es parcialmente anidada: las variantes se comparan dentro de cada control (ej. baseline, istio, kong en C1), aunque ciertas tecnologías como Istio aparecen en múltiples controles. La comparación entre controles no es directa, ya que cada uno representa una clase de mecanismo de protección distinto. Todas las pruebas se ejecutan sobre la misma pila de microservicios y configuración base, permitiendo que las diferencias observadas sean atribuibles principalmente a los controles evaluados.

Esta fase operacionaliza la selección de controles justificada en el capítulo anterior y traslada esa decisión metodológica a un esquema de evaluación empírica con inferencia estadística.

%------------------------------------------------
\section{Entorno experimental}

Los experimentos se ejecutaron en un entorno local controlado. El sistema base es una máquina con procesador AMD Ryzen 5 3600 (6 núcleos, 12 hilos, 3.59 GHz), 16 GB de RAM y Windows 10 Pro (versión 22H2).

La infraestructura se implementó en WSL2 (Ubuntu 20.04) con MicroK8s y containerd, limitando recursos a 4 vCPU, 12 GB RAM y 2 GB swap con localhost forwarding. Aunque inferior al hardware disponible, esta configuración es representativa de entornos de desarrollo profesional. Durante el análisis se consideran las implicaciones de memoria limitada y swap en cargas altas. Cada escenario se ejecutó bajo condiciones controladas, sin cargas concurrentes, manteniendo configuración constante del clúster, asegurando que las diferencias observadas sean atribuibles a los controles de seguridad.
\section{Factores, niveles y variable respuesta}

El factor principal es el tipo de control de seguridad (C), definido como un factor de efectos fijos con cuatro niveles: C1 (API Gateway), C2 (mTLS), C3 (Network Policies) y C4 (Rate Limiting). La selección de estos cuatro controles responde a su relevancia estratégica en arquitecturas de microservicios: conjuntamente forman una secuencia defensiva complementaria que aborda capas distintas del stack (perimetral, transporte, red, aplicación), permitiendo evaluar no solo el impacto aislado de cada mecanismo sino también las interacciones entre protecciones de naturaleza diferente.

Dentro de cada control, se define un segundo factor correspondiente a la variante tecnológica (T), con estructura anidada. Esta estructura refleja realidad práctica: tecnologías como Istio pueden implementar múltiples controles simultáneamente (C1 como Istio Gateway, C2 como mTLS sidecar, C3 mediante virtual services), pero el interés comparativo reside dentro de cada control específico. La variación tecnológica permite discriminar entre el efecto del mecanismo conceptual (presencia/ausencia de autenticación mTLS, por ejemplo) y su implementación particular. Así, C1 compara NGINX Ingress (ingress controller estándar de Kubernetes), Istio Gateway (gateway dentro de malla de servicios completa) e Kong (API Gateway especializado); C2 contrasta baseline sin cifrado (tráfico plaintext) contra Istio mTLS (inyección automática de sidecars de envoy) e Linkerd mTLS (alternativa ligera sin sidecars inyectados); C3 evalúa ausencia de políticas (baseline permisivo), políticas básicas de ingreso y políticas estrictas con énfasis en control de egreso; C4 establece baseline sin limitación de tasa, limitación moderada (1200 requests por minuto) diseñada para permitir carga normal, y limitación estricta (300 rpm) para evaluar degradación controlada bajo presión.

La carga de trabajo (G) no se mantiene constante, sino que se incorpora explícitamente como un factor fijo con cuatro niveles de intensidad: 1, 5, 10 y 20 usuarios virtuales (VUS) ejecutados simultáneamente. La elección de estos niveles sigue una progresión que cubre desde carga baja (1 VUS, evaluando overhead y latencia de seguridad con mínima contención de recursos), carga moderada (5 y 10 VUS, representativa de entornos de desarrollo y testing típicos), hasta carga alta (20 VUS, aproximando límites de infraestructura disponible sin provocar colapso total). Este rango heterogéneo permite observar transiciones de comportamiento: a bajo VUS, los costos fijos de seguridad (handshakes mTLS, validación de políticas) dominan la métrica de latencia; a VUS alto, la contención de recursos y saturación del sistema comienzan a enmascarar diferencias granulares entre controles.

Finalmente, se incorporan réplicas (R) como bloques experimentales organizados aleatoriamente. Las 8 réplicas independientes por celda experimental permiten estimar la variabilidad intrínseca inherente a sistemas distribuidos y mejoran significativamente la robustez de inferencias estadísticas sobre diferencias reales versus ruido experimental. La aleatorización de bloques (no ejecuciones completamente independientes) controla factores de confusión a nivel de sesión: estado residual del clúster, fragmentación de memoria, estados de caché.

En conjunto, el diseño experimental factorial comprende 4 controles, 3 variantes por control y 4 niveles de carga, lo que produce 48 celdas experimentales. Con 8 réplicas por celda, la campaña completa corresponde a 384 ejecuciones brutas, cada una independiente y aislada.

La variable respuesta del estudio está compuesta por un conjunto de seis métricas que caracterizan el comportamiento del sistema bajo seguridad aplicada: latencia media (avg), latencia percentil 95 (p95), throughput (rps), tasa de error, consumo de CPU y consumo de memoria. Estas métricas reflejan tres dimensiones críticas del comportamiento del sistema: (i) rendimiento percibido por usuario (latencias en percentiles), (ii) confiabilidad operativa (tasa de error, throughput sostenido), y (iii) eficiencia de consumo (recursos del clúster). La latencia se mide extremo a extremo desde el cliente (client-side), capturando el impacto completo de los controles de seguridad sin subestimar sobrecostos intermedios; la tasa de error refleja tanto fallos de seguridad intencionados (solicitudes bloqueadas por rate limiting en C4, violaciones de políticas en C3) como fallos técnicos; el throughput y recursos permiten evaluar trade-offs entre seguridad e eficiencia.

\subsection{Clasificación de métricas}

Las métricas se agrupan en tres categorías: (i) rendimiento (latencia y throughput), (ii) confiabilidad (tasa de error), y (iii) eficiencia (consumo de CPU y memoria).


\begin{table}[htbp]
\centering
\caption{Clasificación de variables del estudio}
\label{tab:variables}
\renewcommand{\arraystretch}{1.3}

\begin{tabular}{p{3cm} p{4cm} p{7cm}}
\hline
\textbf{Categoría} & \textbf{Métrica} & \textbf{Descripción} \\
\hline

Rendimiento & Latencia media (avg) & Tiempo de respuesta promedio extremo a extremo bajo carga. \\

Rendimiento & Latencia (p95) & Tiempo de respuesta extremo a extremo medido en el percentil 95. \\

Rendimiento & Tasa de solicitudes (RPS)	 & Cantidad de solicitudes procesadas por segundo. \\

Confiabilidad & Tasa de error & Proporción de solicitudes fallidas sobre el total. \\

Eficiencia & Consumo de CPU & Uso promedio de CPU durante la ejecución, expresado en millicores. \\

Eficiencia & Consumo de memoria & Uso promedio de memoria durante la ejecución, expresado en MiB. \\


\hline
\end{tabular}

\vspace{0.2cm}
{\raggedright \small \textit{Nota.} Clasificación de métricas según su rol en la evaluación del sistema.\par}

\end{table}

%------------------------------------------------
\section{Unidad experimental y réplica}

La unidad experimental es una ejecución completa de un escenario específico caracterizado por una combinación única de control, variante y nivel de carga. Cada ejecución comprende un despliegue limpio, warm-up de 30 s, ventana de carga y cooldown de 15 s.

\begin{itemize}
\item \textbf{Unidad experimental:} una ejecución completa del escenario (control + variante + VUS).
\end{itemize}

Debido a la variabilidad inherente de sistemas distribuidos, se requieren múltiples ejecuciones independientes por escenario para estimar variabilidad y mejorar robustez estadística.

\begin{itemize}
\item \textbf{Número de réplicas:} 8 ejecuciones independientes por escenario, organizadas como bloques aleatorizados.
\item \textbf{Tamaño total del diseño experimental factorial:} 48 celdas experimentales y 384 ejecuciones brutas.
\end{itemize}


%------------------------------------------------

\section{Modelo estadístico}

Se utiliza un modelo factorial con bloques donde el control, variante anidada y nivel de carga son factores de interés, y la réplica actúa como bloque experimental. El modelo para cada métrica es:

\begin{equation}
Y_{ijkl} = \mu + C_i + T_{j(i)} + G_k + B_l + \varepsilon_{ijkl}
\end{equation}

donde $Y_{ijkl}$ corresponde a la variable respuesta observada para el control $i$, la variante $j$ anidada en dicho control, el nivel de carga $k$ y el bloque o réplica $l$; $\mu$ es la media general; $C_i$ representa el efecto del control; $T_{j(i)}$ el efecto de la variante dentro del control; $G_k$ el efecto del nivel de carga; $B_l$ el efecto del bloque; y $\varepsilon_{ijkl}$ representa el error experimental.

Este modelo permite evaluar diferencias significativas entre variantes dentro de cada control y el efecto del nivel de carga bajo un esquema aleatorizado por bloques.

\section{Formulación de hipótesis experimentales}

\begin{table}[htbp]
\centering
\caption{Hipótesis experimentales por control}
\label{tab:hipotesis_controles}
\renewcommand{\arraystretch}{1.4}

\begin{tabular}{p{1.2cm} p{3cm} p{4.5cm} p{4.5cm}}
\hline
\textbf{ID} & \textbf{Factor} & \textbf{Hipótesis nula ($H_0$)} & \textbf{Hipótesis alternativa ($H_1$)} \\
\hline

C1 & API Gateway (NGINX Ingress vs Istio Gateway vs Kong) 
& No existen diferencias estadísticamente significativas en latencia (avg, p95), throughput (rps) y tasa de error entre las configuraciones evaluadas.
& Al menos una configuración presenta diferencias significativas en latencia, throughput o tasa de error. \\

C2 & mTLS (Baseline vs Istio mTLS vs Linkerd mTLS) 
& La implementación de mTLS no produce diferencias significativas en latencia (avg, p95), throughput (rps), tasa de error ni en consumo de recursos (CPU, memoria).
& La implementación de mTLS produce diferencias significativas en al menos una métrica de QoS o consumo de recursos. \\

C3 & Network Policies (Baseline vs Basic vs Strict) & 
La aplicación de políticas de red no genera diferencias significativas en 
latencia (avg, p95), throughput (rps), tasa de error ni en confiabilidad. &
La aplicación de políticas de red genera diferencias significativas en 
al menos una métrica de QoS (latencia, throughput, tasa de error) o confiabilidad. \\


C4 & Rate Limiting (Baseline vs Moderate 1200 rpm vs Strict 300 rpm) 
& La configuración de limitación de tasa no produce diferencias significativas en latencia (avg, p95), throughput (rps), tasa de error ni en consumo de recursos.
& La configuración de limitación de tasa produce diferencias significativas en al menos una métrica de QoS o consumo de recursos. \\

\hline
\end{tabular}

\vspace{0.2cm}
{\raggedright \small \textit{Nota.} Las hipótesis se evalúan mediante análisis de varianza (ANOVA) o pruebas no paramétricas equivalentes, dependiendo del cumplimiento de los supuestos estadísticos.\par}

\end{table}
\FloatBarrier
%------------------------------------------------
